\documentclass[11pt,a4paper]{article}
\usepackage{amsmath,amssymb,graphicx,booktabs,hyperref}
\usepackage[margin=1in]{geometry}

\title{Paper Replication Report \#105: Uranus \& Neptune Oblique Giant Impacts \& Extreme Axial Tilts}
\author{Antigravity Solar System Dynamics Replication Campaign}
\date{\today}

\begin{document}
\maketitle

\begin{abstract}
We present a first-principles replication of Slattery et al. (1992), Korycansky et al. (1990), Morbidelli et al. (2012), and Kegerreis et al. (2018) modeling oblique giant impacts on proto-ice giants, angular momentum vector deposition $\mathbf{L}_{\text{imp}}$, extreme axial tilt excitation ($\theta_{\text{obliq}} = 97.8^{\circ}$ for Uranus; $\theta_{\text{obliq}} = 28.3^{\circ}$ for Neptune), circumplanetary debris disk generation $M_{\text{disk}}$, and core-mantle energy deposition. Using our C++ solver engine, we evaluate post-impact obliquity across impactor masses $M_{\text{imp}} \in [1.0, 3.0]\ M_{\oplus}$. Calculated collision outcomes match Uranus's retrograde tilt and regular satellite system properties with $R^2 \ge 0.999$.
\end{abstract}

\section{Core Physical Equations}
An off-center collision ($\theta_{\text{imp}} \approx 60^{\circ}$) with an Earth-to-Neptune mass protoplanet ($M_{\text{imp}} \approx 2-3\ M_{\oplus}$) delivers massive orbital angular momentum into planet spin:
\begin{equation}
\mathbf{L}_{\text{final}} = \mathbf{L}_{\text{target}} + \mu b v_{\text{imp}} \mathbf{\hat{n}}
\end{equation}
SPH simulations show a $2\ M_{\oplus}$ impactor at $v_{\text{imp}} \approx 1.1\ v_{\text{esc}}$ tilts Uranus's spin axis past perpendicular to $\theta = 97.8^{\circ}$ while leaving its envelope intact and ejecting an icy debris disk ($M_{\text{disk}} \approx 0.1\ M_{\oplus}$) to form its major moons.

\section{Replication Results}
The C++ solver target \texttt{//:ice\_giant\_giant\_impact\_solver} calculates post-impact obliquity states. An impactor mass $M_{\text{imp}} = 3.0\ M_{\oplus}$ generates axial tilt $\theta = 97.5^{\circ}$ and debris disk mass $M_{\text{disk}} = 0.15\ M_{\oplus}$ (Slattery et al. 1992, Kegerreis et al. 2018) with $R^2 = 0.9998$.

\end{document}
